\chapter{Resource limits}
\label{ch:limits}

\begin{aside}
Every TUI has limits: how big a screen it can handle, how
many widgets, how fast it can update. Knowing the limits
prevents the rude surprise.
\end{aside}

\section{Screen size}

\maya{} supports 1x1 through any size your terminal allows.
Very large screens (>500 cells wide) push diff cost above
1 ms. Typical TUIs are under 200x60.

\section{Element count}

A frame with 10,000 elements is fine; layout is O(N) with
small constants. A frame with 100,000 starts to slow.
Virtualise when you hit this.

\section{Signal count}

A few thousand signals are routine. Tens of thousands are
fine. A million might thrash; use stores to bundle related
state.

\section{Effect count}

Each effect costs a function call on every relevant signal
change. Hundreds are fine. Thousands of effects, each with
many subscribers, can dominate the frame.

\section{Memory}

A typical \maya{} app uses 10--50 MB resident. Heavy data
(large document loaded) drives this. Use lazy loading and
virtualisation.

\section{Wire bandwidth}

A frame's ANSI output is typically 100 bytes to 5 KB. Full
redraws hit 50 KB or more. On slow links, frame budget is
the bottleneck.

\section{File descriptors}

\maya{} uses 2--5 fds at runtime (stdin, stdout, log file).
Adding watchers, sockets, child PTYs grows this. Default
ulimit is 1024 on Linux; pay attention if you embed many
PTYs.

\section{Recap}

\begin{recap}
\begin{itemize}[leftmargin=*]
  \item Screen, elements, signals, effects, memory,
    wire, fds.
  \item Each has a sweet spot for typical TUIs.
  \item Know yours; design within.
\end{itemize}
\end{recap}

\section*{Workshop}

\begin{enumerate}[leftmargin=*]
  \item Measure your app's resource usage.
  \item Identify the limit you'll hit first.
  \item Plan how to scale past it.
\end{enumerate}
